SR Resolution and signal shape (after some GNN cut)\\

We study the resonant signal shape without the jpsi contrained B mass
as proxy for the non-resonant signal shape.\\
1/ signal shape: double-sided Crystal Ball, i.e. with single mean and two independent
left and right sigmas and tail parameters.\\
insert plots of signal shapes from MC:\\
a/ res vs non-res without q2 cuts\\
b/ with q2 cut: q2 cut modifies the low tail as can also be seen by the 2D distribution plot in (previous section)\\
c/ fixed the shape such that we can extract the non-res signal shape from the resonant fit by correcting for the q2 effect.\\

Want to check if signal shape is compatible between high- and low- q2 channels. Comparing the fit shape (to MC)\\
Good agreement of signal shape between channels -> good as we want to take advantage of high stats in Res channel to fix NR channel signal shape\\
Due to correlation of B mass and q2, a cut on q2 corresponds to clipping the high and low mass tails of the resonant channel -> can distort shape.\\
This makes the mass distribution take a less easily modelled form, and so need to account for this if we want to do simultaneous fit\\

Background modelling strategy is same as in the jpsi-contrained B mass
fits: only the B+ needs to be rethought as the GNN could cut its
yields and it might feed in the background due to the resolutions.\\
Possibly this could be included as a systematic effect.\\
To be checked: GNN efficiency on B+ MC and SS jpsi-constrained and
not jpsi-constrained fits after GNN.\\

Also the leakage from the resonant channel into the low q2 region
needs to be studied. Not necessarily a separate component in the fit
as it merges with a exponential shape, however a systematic can be
added for this.\\

Insert resonance fit results with plot and table. And then TOYS!\\

Eta-dependent signal fit: on resonant only fits with toys. Assess significance
and stability? (think about what figure of merit we want to use to
decide this. In general as this adds complexity given the already
rather complex fits, as we want to have a simultaneus fit with the res and
non-res, one might prefer to go for an eta-independent case if the difference
in performance is negligible.)\\

Simultaneous fits: non-res and res and SS

